%---------------------------------------------------------------------
\begin{abstract}
Hangboarding---isometric dead-hang training from purpose-built holds---is the best-studied supplemental training method in rock climbing. Ten intervention trials (nine randomised controlled trials and one controlled study, combined $n \approx 230$) and one large retrospective cohort ($n = 527$, published 2012--2025), together with one systematic review and meta-analysis (Stien et al.\ 2023, $n = 110$ from 5 trials), support three principal conclusions: (1) adding 2--3 hangboard sessions per week to regular climbing consistently improves maximal finger strength (MFS) and grip endurance within 4--10 weeks in intermediate-to-advanced male climbers (meta-analytic SDM~= 0.41 for MFS, 95\%~CI 0.03--0.80; SDM~= 0.57 for overall climbing-specific test performance, 95\%~CI 0.24--0.91); (2) $\geq\!80\%$-MFS protocols maximise MFS, 60--80\% MFS repeater protocols maximise stamina, and 80\% MFS improves both simultaneously; (3) daily low-intensity hangs (${\sim}40\%$ MFS, ``Abrahangs'') produced a positive but non-significant trend versus climbing-only controls ($+2.5\%$, $p = 0.12$) while twice-weekly max hangs were significant ($+3.2\%$, $p < 0.001$); combining both modalities was additive ($+5.8\%$). These Abrahangs findings are from a retrospective observational cohort and are hypothesis-generating only. The literature is uniformly limited by small samples (PEDro scores 5--7/10), male-advanced cohorts, short follow-up ($\leq10$~weeks), and surrogate (dynamometric) outcomes. All effect sizes are on surrogate measures; no trial has demonstrated statistically significant climbing-grade improvement attributable to hangboarding alone.
\end{abstract}
